%% notes-tex/031-yeast-chemical-space/sections/4-open.tex

\section{What is left, and the question for a collaborator\secstatus{todo}}
\label{sec:open}

Three openings survive the published work, and they are narrower and more specific than the
original framing.

\notesec{Stereochemical resolution, which nobody reports}
A rule-based expansion operates on SMILES, and 621 of the 894 Yeast9 structures name a set of
stereoisomers rather than a molecule. A species that is stereochemically ambiguous cannot
carry an atom mapping that means anything, cannot be assigned a unique thermodynamic value,
and collapses L and D forms onto one key. Resolving the 1{,}378 species to single molecules is
a prerequisite for a physical accounting, it is measured here as a gap, and no paper in the
mirror reports it for any yeast reconstruction. This is the cleanest opening.

\notesec{Lipids, where both the model and the expansion are weakest}
Lipids dominate every count. Of Yeast9's 1{,}378 species, \textbf{281 are acyl-resolved or
pooled lipid classes and not one carries a structure}. Of the 14{,}882 metabolites the
published comparison found missing, 14{,}310 are lipids. And the published expansion reaches
92 percent of the whole metabolome but only 75 percent of the non-lipid subset, which means
the whole-metabolome figure is carried by lipids enumerated combinatorially from acyl chains
rather than by resolved chemistry. Whether that enumeration is a physical accounting or a
parameterization is a fair question to put to a specialist.

\notesec{The mass fraction, which is simply missing}
Nothing in the mirrored literature states what share of yeast dry mass is small molecules.
Until that number exists the exercise has no denominator: an accounting of every chemical
species is a different undertaking if the pool is 3 percent of the cell than if it is 15.

\notesec{Does retrobiosynthesis beat constraint-based gap filling here}
Worth stating precisely, because the two do different jobs. Constraint-based gap filling adds
reactions from a universal database until a model carries flux to a required biomass
component, so it is driven by a phenotype and answers with the smallest set of reactions that
restores it. It never proposes a molecule that is not already in its database. A
retrobiosynthetic expansion applies reaction rules to structures and generates species that
were in no database, so it can propose the intermediate itself. For enumerating the space of
chemical intermediates, which is the stated goal, only the second can work. The measurement in
Section~\ref{sec:inventory} supports that from the other direction: reachability on the
species graph saturates at 99.7 percent in four rounds and therefore bounds nothing, so the
bound has to come from chemistry.

\textit{Hypothesis (untested):} the residual the published expansion reports, 267 non-lipid
metabolites that could not be connected with enzyme-annotated reactions, would shrink if
uncatalyzed chemistry were admitted on the same footing as enzymatic rules, since spontaneous
and metal-catalyzed steps are exactly what an enzyme-annotated rule set excludes. The
published work carries 213 spontaneous rules against 21{,}921 enzymatic ones, which is a small
allowance for a cell held at 30 degrees for many hours.

\notesec{The concrete ask}
Put to a retrobiosynthesis specialist, the question is not whether yeast metabolism can be
expanded, because it has been. It is three narrower things. Can a rule set operate on
stereochemically resolved species, and what does it cost to resolve 1{,}378 of them first.
Does admitting uncatalyzed and metal-catalyzed chemistry close the 267-metabolite residual, or
does it inflate the space without bound. And is there a representation of acyl-resolved lipids
that is a physical accounting rather than a combinatorial enumeration.

The deliverable that would make the conversation concrete is a conversion of Yeast9 into a
reaction format a retrobiosynthesis tool consumes, which means per-reaction substrate and
product structures with an atom mapping where one can be derived. Section~\ref{sec:inventory}
measures the ceiling on that conversion as released: \textbf{2{,}156 of 4{,}131 reactions},
52 percent, have a structure for every participant today. That conversion is a small piece of
work and it is the natural first step, but it is a side quest and it does not block the
modeling phase of experiment 031.
